\documentclass[12pt]{article}
\usepackage[T1]{fontenc}
\usepackage[utf8]{inputenc}
\usepackage{lmodern}
\usepackage{textcomp}
\usepackage{enumitem}
\usepackage{geometry}
\geometry{margin=1in}
\begin{document}
\begin{center}
Answer Key - Psychology - Graduate Level Assessment 3 \
\small (Answers are ordered exactly as in the question paper.)
\end{center}

\section*{Multiple Choice}
\begin{enumerate}[label=\arabic*.]
\item \textbf{Answer:} A
\\ \textit{Options:} A) Content validity, criterion validity, construct validity, face validity; B) Test-retest, internal consistency, inter-rater, content validity; C) Alternate form, inter-rater, construct validity, criterion-related validity; D) Test-retest, split-half, inter-rater, test validity; E) Test-retest, inter-rater, intra-rater, predictive validity
\\ \textit{Note:} The correct answer identifies four essential types of validity-content validity, criterion validity, construct validity, and face validity-that collectively assess the accuracy and appropriateness of psychological measurements, ensuring they effectively evaluate the constructs they are intended to measure.
\item \textbf{Answer:} A
\\ \textit{Options:} A) a learned skill that improves with practice; B) a definitive decision made based on facts; C) an emotional state that varies frequently; D) a set of behaviors that are consistent across different situations; E) a moral judgment that defines right or wrong actions
\\ \textit{Note:} An attitude is considered a learned skill because it is shaped by experiences and social influences, and like any skill, it can be developed and refined through deliberate practice and exposure to new information.
\item \textbf{Answer:} A
\\ \textit{Options:} A) identifying the causes and solutions of the client's presenting problems; B) establishing a hierarchy of authority to enable effective decision making; C) ensuring the consultee adheres strictly to a predetermined action plan; D) proposing multiple alternative solutions for the consultee to choose from; E) identifying the strengths and weaknesses of the consultee's current approach
\\ \textit{Note:} In Caplan's model of consultee-centered case consultation, the consultant focuses on helping the consultee understand and address the underlying causes and potential solutions to the client's presenting issues, emphasizing a collaborative approach to problem-solving.
\item \textbf{Answer:} B
\\ \textit{Options:} A) a verbal praise or encouragement.; B) a behavior that occurs frequently.; C) a specific reward given after the desired behavior.; D) a stimulus that is artificially created.; E) a generalized conditioned reinforcer.
\\ \textit{Note:} The Premack Principle states that a more probable behavior can be used as a reinforcer for a less probable behavior, meaning that a behavior that occurs frequently can effectively reinforce a less frequent behavior to increase its occurrence.
\item \textbf{Answer:} B
\\ \textit{Options:} A) are more creative but worse at decision-making.; B) are more creative and better at decision-making.; C) make better decisions but are less creative overall.; D) are less creative and less productive overall.; E) are more productive but worse at decision-making.
\\ \textit{Note:} Heterogeneous work groups bring together diverse perspectives and experiences, which fosters creativity and enhances decision-making by encouraging innovative solutions and reducing groupthink.
\end{enumerate}

\section*{True / False}
\begin{enumerate}[label=\arabic*.]
\setcounter{enumi}{5}
\item \textbf{Answer:} False
\\ \textit{Options:} A) True; B) False
\\ \textit{Note:} The variance of the sample observations 2, 5, 7, 9, and 12 is calculated to be 11.2, not 8.5, as variance is determined by the average of the squared differences from the mean.
\item \textbf{Answer:} True
\\ \textit{Options:} A) True; B) False
\\ \textit{Note:} Systematic desensitization primarily utilizes gradual exposure to anxiety-provoking stimuli along with relaxation techniques such as deep breathing or progressive muscle relaxation, rather than hypnosis, to help clients manage their fears.
\item \textbf{Answer:} False
\\ \textit{Options:} A) True; B) False
\\ \textit{Note:} Quasi-experimental designs do not involve random assignment of participants to groups; instead, they rely on existing groups or non-randomized methods to assign participants, which distinguishes them from true experimental designs.
\item \textbf{Answer:} True
\\ \textit{Options:} A) True; B) False
\\ \textit{Note:} In Gestalt therapy, transference is considered a projection of the client's internal conflicts and unresolved issues onto the therapist, thus reflecting their subjective experiences and fantasies rather than a direct representation of the therapist.
\item \textbf{Answer:} False
\\ \textit{Options:} A) True; B) False
\\ \textit{Note:} The localization of sounds is primarily influenced by spatial cues such as interaural time differences and interaural level differences, rather than pitch or dynamic factors.
\end{enumerate}

\section*{Long Form Response}
\begin{enumerate}[label=\arabic*.]
\setcounter{enumi}{10}
\item \textbf{Answer:} Binocular vision significantly enhances depth perception through the mechanisms of binocular disparity and the interpretation of double images. Each eye captures a slightly different image due to their spatial separation, creating binocular disparity, which the brain processes to gauge depth. This disparity provides critical information about the relative distances of objects in the visual field. Furthermore, when objects are closer than the point of convergence, they produce double images, prompting the brain to interpret these images as indicating proximity. Together, these processes yield a more nuanced understanding of spatial relationships, allowing individuals to navigate their environment effectively.
\\ \textit{Note:} correct because it accurately describes how binocular disparity arises from the different perspectives of each eye, enabling the brain to interpret depth and distance, while also addressing the phenomenon of double images for close objects, which further aids in depth perception.
\item \textbf{Answer:} Maria's fear of success can be attributed to several psychological factors, including a fear of increased expectations, pressure to perform, and the potential for isolation from peers. This fear may lead her to favor community colleges, particularly local ones, as they often represent a more manageable and less intimidating step in her academic journey. By choosing a familiar, nearby institution, she may seek to mitigate the anxiety associated with larger, prestigious universities, which could evoke feelings of inadequacy and fear of failure. Conversely, distant community colleges might symbolize a leap towards greater achievement, triggering her fear of success and leading her to avoid such opportunities altogether. Ultimately, her decisions regarding college applications will likely reflect an underlying conflict between aspiration and anxiety, shaping her educational trajectory.
\\ \textit{Note:} correct because it identifies specific psychological factors, such as fear of increased expectations and potential isolation, that contribute to Maria's fear of success, and it logically connects these factors to her preference for the less intimidating environment of community colleges.
\end{enumerate}

\vfill
\noindent\textit{End of Answer Key}
\end{document}